\documentclass[10pt]{article}

% ---------- Document Identity (update these only) ----------
\newcommand{\DocStage}{}
\newcommand{\DocTitle}{Your Road to Tier One SOF}
\newcommand{\ProgramName}{182-Week Civilian-to-Selection Readiness Pipeline}
\newcommand{\ProgramSubtitle}{}

% ---------- Page ----------
\usepackage[legalpaper,left=0.5in,right=0.5in,top=0.6in,bottom=0.45in]{geometry}

% ---------- Typography ----------
\usepackage[T1]{fontenc}
\usepackage[utf8]{inputenc}
\usepackage{libertinus}
\usepackage{microtype}
\usepackage{parskip}
\usepackage{ragged2e}
\usepackage{textcomp}
\AtBeginDocument{\color{black}}
\AtBeginDocument{\pagecolor{white}}
\AtBeginDocument{\justifying}

% ---------- Landscape pages ----------
\usepackage{pdflscape}

% ---------- Color ----------
\usepackage[table]{xcolor}
\definecolor{StageBlue}{HTML}{1F4E79}
\definecolor{StageBlueDark}{HTML}{163A5A}
\definecolor{LightGray}{HTML}{F2F4F7}
\definecolor{MidGray}{HTML}{D9DEE5}
\definecolor{RuleGray}{HTML}{C7CED8}
\definecolor{HeaderInk}{HTML}{000000}
\definecolor{Ink}{HTML}{000000}

% ---------- Utility ----------
\usepackage{enumitem}
\sloppy
\setlength{\emergencystretch}{2em}

% ---------- Boxes / UI ----------
\usepackage[most]{tcolorbox}

\tcbset{
  charterTitle/.style={
    enhanced,
    colback=StageBlue!4,
    colframe=StageBlue,
    boxrule=1.25pt,
    arc=3pt,
    left=16pt,right=16pt,top=14pt,bottom=14pt,
    borderline north={3.2pt}{0pt}{StageBlue},
    borderline west={4pt}{0pt}{StageBlue},
    borderline south={1.25pt}{0pt}{StageBlue},
    drop shadow={black!25!white},
  },
  charterRule/.style={
    enhanced,
    colback=LightGray,
    colframe=RuleGray,
    boxrule=0.85pt,
    arc=3pt,
    left=12pt,right=12pt,top=10pt,bottom=10pt,
    borderline west={3pt}{0pt}{StageBlue},
  },
  charterBadge/.style={
    enhanced,
    colback=StageBlue,
    coltext=white,
    colframe=StageBlueDark,
    boxrule=0.6pt,
    arc=2pt,
    left=6pt,right=6pt,top=3pt,bottom=3pt,
  }
}
\newtcbox{\Badge}{nobeforeafter,charterBadge}

% ---------- Lists ----------
\setlist[itemize]{leftmargin=1.5em, itemsep=0.25em, topsep=0.25em}
\setlist[enumerate]{leftmargin=1.8em, itemsep=0.25em, topsep=0.25em}

% ---------- TOC ----------
\usepackage{tocloft}
\setcounter{tocdepth}{3}
\setlength{\cftbeforesecskip}{3pt}
\setlength{\cftbeforesubsecskip}{2pt}
\setlength{\cftbeforesubsubsecskip}{0pt}
\renewcommand{\cftsecfont}{\bfseries}
\renewcommand{\cftsecpagefont}{\bfseries}
\renewcommand{\cftsubsecfont}{\normalfont}
\renewcommand{\cftsubsecpagefont}{\normalfont}
\renewcommand{\cftsubsubsecfont}{\normalfont}
\renewcommand{\cftsubsubsecpagefont}{\normalfont}
\setlength{\cftsecindent}{0pt}
\setlength{\cftsubsecindent}{1.25em}
\setlength{\cftsubsubsecindent}{2.8em}
\setlength{\cftsecnumwidth}{2.1em}
\setlength{\cftsubsecnumwidth}{2.8em}
\setlength{\cftsubsubsecnumwidth}{3.6em}

% Remove the built-in TOC heading so only the custom heading is shown
\makeatletter
\renewcommand{\tableofcontents}{\@starttoc{toc}}
\makeatother

% ---------- Header/Footer: page number only ----------
\usepackage{fancyhdr}
\pagestyle{fancy}
\fancyhf{}
\renewcommand{\headrulewidth}{0pt}
\renewcommand{\footrulewidth}{0pt}
\fancyfoot[C]{\thepage}

% ---------- Section formatting ----------
\usepackage{titlesec}

\titleformat{\section}
  {\Large\bfseries\color{Ink}}
  {\thesection}{0.6em}{}
  [\vspace{0.25em}\textcolor{StageBlue}{\rule{\linewidth}{0.8pt}}]

\titleformat{\subsection}
  {\large\bfseries\color{Ink}}
  {\thesubsection}{0.6em}{}

\titleformat{\subsubsection}
  {\normalsize\bfseries\color{Ink}}
  {\thesubsubsection}{0.6em}{}

\titlespacing*{\section}{0pt}{0.95em}{0.35em}
\titlespacing*{\subsection}{0pt}{0.65em}{0.25em}
\titlespacing*{\subsubsection}{0pt}{0.55em}{0.2em}

% ---------- Table engine ----------
\usepackage{tabularray}
\UseTblrLibrary{booktabs}

% ---------- tabularray: GLOBAL table treatment ----------
\SetTblrInner{
  cells = {fg=black, bg=white, valign=t},
}

% ---------- Header helpers ----------
\newcommand{\Hdr}[1]{\shortstack[c]{\strut #1\strut}}

% ---------- Lines ----------
\newcommand{\TitleLine}{\textcolor{StageBlue}{\rule{0.98\linewidth}{0.95pt}}}
\newcommand{\SubTitleLine}{\textcolor{RuleGray}{\rule{0.98\linewidth}{0.65pt}}}

% ---------- Continued on next page ----------
\newcommand{\ContinuedOnNextPageInline}{%
  \par
  \vfill
  \noindent\textcolor{RuleGray}{\rule{\linewidth}{1pt}}\vspace{-2pt}
  \noindent\hfill\textit{Continued on next page}\par\vspace{-10pt}
  \noindent\textcolor{RuleGray}{\rule{\linewidth}{1pt}}
  \clearpage
}

% ---------- PDF hyperlinks ----------
\usepackage[hidelinks]{hyperref}
\hypersetup{
  colorlinks=false,
  pdfborder={0 0 0},
  linktoc=all,
  pdftitle={\DocTitle},
  pdfauthor={\ProgramName}
}

% =========================================================
\begin{document}

\section*{7.02 — Hormonal Support and Adaptogens — OTC}
\phantomsection
\addcontentsline{toc}{section}{7.02 — Hormonal Support and Adaptogens — OTC}

\subsection*{7.02.1 Purpose \& Scope}
\phantomsection
\addcontentsline{toc}{subsection}{7.02.1 Purpose \& Scope}

7.02 exists to define education-only boundaries for OTC adaptogens and OTC “hormonal support” marketing categories, with strict clinician escalation posture for endocrine concerns. It is not a performance lever and is never required for any phase. It cannot be used to justify gate outcomes or to override safety, recovery, and validity posture.

\subsection*{7.02.2 Governance and Safety Boundaries}
\phantomsection
\addcontentsline{toc}{subsection}{7.02.2 Governance and Safety Boundaries}

\begin{itemize}
\item OTC-only for any item discussed here.
\item Endocrine concerns are not self-managed: seek licensed clinician evaluation when trigger symptoms or abnormal labs occur.
\item Prohibited content and behavior:
  \begin{itemize}
  \item no advice on anabolic agents, controlled substances, SARMs, peptides, research chemicals, or prescription-only hormones,
  \item no dosing/sourcing/cycling/procurement guidance for prescription hormones,
  \item no “post-cycle” or hormone manipulation strategies.
  \end{itemize}
\item Validity protection:
  \begin{itemize}
  \item do not introduce new supplement variables near Deload+Test windows or gate attempts,
  \item do not change supplement variables to “rescue” a poor test week or low-compliance window.
  \end{itemize}
\item Safety-first: any adverse reaction, red-flag symptom, or unusual physiologic response triggers \textbf{STOP} \textbf{NOW} posture for that day’s testing/exposure and clinician escalation when indicated.
\end{itemize}

\subsection*{7.02.3 Tier Doctrine Applied to 7.02}
\phantomsection
\addcontentsline{toc}{subsection}{7.02.3 Tier Doctrine Applied to 7.02}

Tiering for 7.02 represents the degree of education and documentation posture, not a required acquisition pathway.

\begin{itemize}
\item Tier 0: education-only boundary awareness; no product required; no execution dependency.
\item Tier 1: optional single-variable trial posture (one item at a time), conservative dosing range context, and strict documentation; still never required.
\item Tier 2: optional expanded education depth and tighter documentation controls to reduce novelty and interaction risk; still never required.
\item Tier 3: optional overbuilt documentation and risk-control posture (redundancy in tracking, conservative stop rules, and clinician coordination); still never required.
\end{itemize}

In every phase and every tier, 7.02 remains non-required and cannot be treated as a gate dependency.

\subsection*{7.02.4 Selection Criteria \& Fit Checks}
\phantomsection
\addcontentsline{toc}{subsection}{7.02.4 Selection Criteria \& Fit Checks}

\begin{itemize}
\item Eligibility filter (non-diagnostic):
  \begin{itemize}
  \item if the trainee is already using prescription medications (especially thyroid, psychiatric, anticoagulants, blood pressure), treat any adaptogen use as clinician-led discussion only.
  \item if the trainee has a history of severe anxiety/panic, insomnia vulnerability, or adverse reactions to stimulants, avoid activating adaptogens and treat as omit-by-default.
  \end{itemize}
\item Product type selection (generic):
  \begin{itemize}
  \item prefer single-ingredient items to reduce interaction uncertainty.
  \item avoid multi-ingredient “hormone boosters” and proprietary blends; they increase novelty risk and make adverse-event attribution impossible.
  \end{itemize}
\item Stability requirement:
  \begin{itemize}
  \item one variable at a time; do not change dose and item simultaneously.
  \item do not initiate or change within Deload+Test windows.
  \end{itemize}
\item Fit check outcomes (operational):
  \begin{itemize}
  \item “Fit” means no adverse effects, no sleep disruption, no anxiety escalation, no GI distress that affects training execution, and no changes that contaminate gates.
  \item “Not fit” means any adverse event trend or any need to chase dose/stack multiple items to “feel it,” which is prohibited.
  \end{itemize}
\end{itemize}

\subsection*{7.02.5 Setup, Safety, and Anchoring}
\phantomsection
\addcontentsline{toc}{subsection}{7.02.5 Setup, Safety, and Anchoring}

\begin{itemize}
\item Setup posture (non-prescriptive):
  \begin{itemize}
  \item choose one optional item at a time, keep timing stable if used, and document as a controlled variable.
  \item do not use adaptogens as a pre-test acute performance enhancer; this contaminates comparability and violates “no novelty near gate” posture.
  \end{itemize}
\item Anchoring to governance:
  \begin{itemize}
  \item treat 7.02 as never required; omit is always acceptable.
  \item any endocrine concern triggers clinician evaluation; do not self-diagnose.
  \end{itemize}
\item Stop posture for adverse events:
  \begin{itemize}
  \item discontinue the optional item if adverse effects occur (insomnia, anxiety, palpitations, rash, GI distress, unusual dizziness).
  \item if the adverse event is severe or includes red-flag symptoms, \textbf{STOP} \textbf{NOW} posture applies and escalation is required per governance.
  \end{itemize}
\end{itemize}

\subsection*{7.02.6 Maintenance, Replacement, and Storage}
\phantomsection
\addcontentsline{toc}{subsection}{7.02.6 Maintenance, Replacement, and Storage}

\begin{itemize}
\item Storage:
  \begin{itemize}
  \item store in a cool, dry location; avoid heat/humidity exposure that degrades stability.
  \item keep containers closed; prevent moisture contamination.
  \end{itemize}
\item Replacement:
  \begin{itemize}
  \item respect expiration dates; discard expired products.
  \item discontinue and discard if odor/appearance changes or contamination is suspected.
  \end{itemize}
\item Documentation:
  \begin{itemize}
  \item maintain a simple log record (item type, dose, timing, start/stop date, adverse events).
  \item keep label photo and lot information where available, to support audit and contamination risk mitigation.
  \end{itemize}
\end{itemize}

\subsection*{7.02.7 Substitution Rules}
\phantomsection
\addcontentsline{toc}{subsection}{7.02.7 Substitution Rules}

\begin{itemize}
\item Primary substitution for any optional 7.02 item is omit. This category is never required.
\item If the intent is “stress support,” substitute with governance-aligned controls:
  \begin{itemize}
  \item downshift training intensity/density,
  \item protect sleep opportunity,
  \item apply Stage 1.F and Stage 3 reslot discipline rather than forcing gate attempts.
  \end{itemize}
\item If endocrine symptoms appear, substitute with clinician evaluation:
  \begin{itemize}
  \item stop self-directed experimentation,
  \item provide clinician with symptom log and medication/supplement disclosure,
  \item do not resume until cleared.
  \end{itemize}
\item Validity substitution rule:
  \begin{itemize}
  \item do not introduce or change 7.02 variables inside protected windows; if a change occurs, treat gate attempts as validity-sensitive and reslot rather than claiming gate credit.
  \end{itemize}
\end{itemize}

\subsection*{7.02.8 Phase Integration Map}
\phantomsection
\addcontentsline{toc}{subsection}{7.02.8 Phase Integration Map}

\begingroup
\scriptsize
\setlength{\tabcolsep}{2pt}
\renewcommand{\arraystretch}{1.12}

\begin{longtblr}{
  width=\linewidth,
  rowhead=1,
  colspec = {
    Q[l,wd=1.05in]
    Q[l,wd=1.70in]
    X[l,2.25]
    X[l,3.75]
  },
  hlines,
  vlines,
  cells = {hsep=6pt, vsep=6pt, halign=l, valign=t},
  row{odd} = {bg=white},
  row{even} = {bg=LightGray},
  row{1} = {bg=StageBlue!15, fg=black, font=\bfseries, halign=l, valign=t},
}
\Hdr{Phase} &
\Hdr{Required Tier (from Stage 6)} &
\Hdr{What Changes / Becomes Mandatory} &
\Hdr{Notes} \\
Phase 00 & T0 & None & Never required. Strongly discourage novelty near Deload+Test windows; focus is compliance, sleep, and low-risk tolerance. \\
Phase 01 & T0 & None & Never required. If used at all, stabilize well before protected windows and document as context only. \\
Phase 02 & T0 & None & Never required. If \textbf{STR} consequence increases, do not use 7.02 to mask recovery debt or to justify escalation. \\
Phase 03 & T0 & None & Never required. Treat any change as a controlled variable; no introductions near gates. \\
Phase 04 & T0 & None & Never required. Aquatic phase relevance does not change 7.02 posture. \\
Phase 05 & T0 & None & Never required. Longer endurance windows increase the importance of sleep and hydration; do not attempt to “supplement around” instability. \\
Phase 06 & T0 & None & Never required. Higher consequence gates demand lower variance; avoid novel supplements near protected windows. \\
Phase 07 & T0 & None & Never required. Hybrid density increases drift risk; governance response is reslot/downshift, not supplement escalation. \\
Phase 08 & T0 & None & Never required. Durability constraints dominate; treat 7.02 as omit-by-default. \\
Phase 09 & T0 & None & Never required. \textbf{CAL} volume is controlled through programming and recovery posture, not adaptogen reliance. \\
Phase 10 & T0 & None & Never required. Underwater governance is unrelated; 7.02 remains education-only. \\
Phase 11 & T0 & None & Never required. Late-phase validity demands discourage novelty; omit remains preferred. \\
Phase 12 & T0 & None & Never required. Verification blocks require low variance; do not add or change supplements inside protected windows. \\
Phase 13 & T0 & None & Never required. Consolidation demands repeatability; omit remains preferred. \\
\end{longtblr}

\endgroup

Phase Integration Map Notes:

\begin{itemize}
\item Stage 6 represents 7.02 as education engagement posture across phases. This overview treats required tier as T0 for every phase because 7.02 is never required; any education engagement remains optional.
\end{itemize}

\subsection*{7.02.9 Risks, Contraindications, and Red Flags}
\phantomsection
\addcontentsline{toc}{subsection}{7.02.9 Risks, Contraindications, and Red Flags}

High-level risks (education-only; non-diagnostic):

\begin{itemize}
\item Interaction risk:
  \begin{itemize}
  \item sedatives/sleep medications, antidepressants, stimulants, thyroid medications, anticoagulants, and blood pressure medications can interact with OTC botanicals.
  \item default posture: if interaction risk is plausible or uncertain, omit and seek clinician input.
  \end{itemize}
\item Sleep disruption and anxiety amplification:
  \begin{itemize}
  \item common failure mode is worsening sleep quality, anxiety, agitation, or palpitations, which degrades recovery and contaminates gate evidence.
  \end{itemize}
\item GI intolerance and dehydration risk:
  \begin{itemize}
  \item nausea, diarrhea, appetite disruption can reduce training feasibility and increase safety risk.
  \end{itemize}
\item Contamination and mislabeling risk:
  \begin{itemize}
  \item OTC supplements can have variable purity; treat this as a validity and safety risk, especially when selection-readiness contexts may require clean compliance posture.
  \end{itemize}
\item Endocrine red flags (trigger clinician evaluation; do not self-manage):
  \begin{itemize}
  \item persistent low libido or sexual dysfunction,
  \item unexplained fatigue that does not respond to sleep normalization,
  \item unexplained weight change with temperature intolerance,
  \item persistent mood changes, new anxiety/panic, or sleep collapse,
  \item new breast tenderness/swelling, unexplained bruising, or edema.
  \end{itemize}
\item Common lab abnormalities that should trigger clinician evaluation (non-diagnostic examples):
  \begin{itemize}
  \item abnormal TSH/free T4,
  \item abnormal total/free testosterone with abnormal LH/FSH,
  \item abnormal prolactin,
  \item unexplained elevated liver enzymes,
  \item abnormal fasting glucose/A1c,
  \item abnormal hematocrit/hemoglobin changes,
  \item abnormal lipid profile changes not explained by diet/execution context.
  \end{itemize}
\end{itemize}

Stop posture:

\begin{itemize}
\item If severe symptoms occur (chest pain/pressure, syncope/near-syncope, severe shortness of breath, new neurologic symptoms), \textbf{STOP} \textbf{NOW} and escalation posture applies per governance. Do not attribute symptoms to supplements as a “normal reaction.”
\end{itemize}

Requires Stage 8 review triggers (governance drift):

\begin{itemize}
\item Any downstream artifact that attempts to treat 7.02 as required, or introduces prohibited prescription hormone guidance, requires Stage 8 review and renders the artifact non-deployable until corrected.
\end{itemize}

\subsection*{7.02.10 Compliance Checklist}
\phantomsection
\addcontentsline{toc}{subsection}{7.02.10 Compliance Checklist}

\begin{itemize}
\item \textbf{PASS}/\textbf{FAIL}: Category is explicitly OTC-only and education-only; no prescription hormone guidance appears.
\item \textbf{PASS}/\textbf{FAIL}: No dosing/sourcing/cycling/procurement guidance for prescription hormones appears; no prohibited substances are discussed.
\item \textbf{PASS}/\textbf{FAIL}: 7.02 is explicitly never required for any phase; no gate dependency posture is implied.
\item \textbf{PASS}/\textbf{FAIL}: Clinician escalation triggers include endocrine symptom examples and common lab abnormality examples without diagnostic claims.
\item \textbf{PASS}/\textbf{FAIL}: Validity contamination controls are explicit (no novelty near Deload+Test windows; controlled variable documentation).
\item \textbf{PASS}/\textbf{FAIL}: Water governance and broader safety boundaries are not contradicted (no breath-hold/underwater instruction; no unsafe claims).
\end{itemize}

\end{document}